\begin{theorem}
$\bpp \subset \ppoly$.
\end{theorem}

\begin{proof} Пусть $A \in \bpp$. Рассмотрим алгоритм $V$ такой, дающий для слов длины $n$ верный ответ с вероятностью $1 - \eps$, $\eps < \frac{1}{2^n}$. Тогда вероятность того, что алгоритм даст неверный ответ хотя бы для одного $x$ равна $2^n \cdot \eps < 1$ (всего слов $2^n$, вероятность ошибиться - $\varepsilon$), а значит вероятность, что на всех $x$ алгоритм отработает корректно ненулевая. Отсюда получаем, что $\exists r_0 \ \brackets{\forall x: |x| = n \ V(x, r_0) = 1 \Longleftrightarrow x \in A}$. Поскольку $V$ работает за полиномиальное время, стандартное преобразование в схему с <<запаянным>> $r_0$ даст схему полиномиального размера, распознающую все $x$ длины $n$.

\end{proof}